\documentclass[preview]{standalone}

\usepackage{amsmath}
\usepackage{amssymb}
\usepackage{stellar}
\usepackage{definitions}

\begin{document}

\id{psicologia-sensazione-percezione}
\genpage

\section{Sensazione e percezione}

\begin{snippet}{sensazione-percezione-introduzione}
    Sensazione e percezione forniscono il materiale di base per i processi
    cognitivi. Tuttavia, i sensi non registrano semplicemente gli stimoli: il
    materiale sensoriale viene elaborato in tempo reale da processi cognitivi
    che generano interpretazioni del mondo a partire dagli stimoli esterni e
    dalle conoscenze precedenti.
\end{snippet}

\begin{snippetdefinition}{percezione-definition}{Percezione}
    La \emph{percezione} è il processo generale attraverso cui un organismo
    individua oggetti ed eventi nell'ambiente, attribuisce loro un senso, li
    comprende, li riconosce e si prepara ad agire rispetto a essi.
\end{snippetdefinition}

\begin{snippetdefinition}{percetto-definition}{Percetto}
    Un \emph{percetto} è ciò che viene percepito, cioè il risultato fenomenico
    di cui il soggetto fa esperienza attraverso il processo percettivo.
\end{snippetdefinition}

\begin{snippet}{fasi-processo-percettivo}
    Il processo percettivo può essere distinto in tre fasi:
    \begin{enumerate}
        \item \emph{sensazione};
        \item \emph{organizzazione percettiva};
        \item \emph{identificazione e riconoscimento degli oggetti}.
    \end{enumerate}
\end{snippet}

\includesnpt[width=62\%,src=/snippet/static-psychology/perception-process-stages.jpg]{centered-img}

\begin{snippetdefinition}{sensazione-definition}{Sensazione}
    La \emph{sensazione} è l'impressione soggettiva, immediata e semplice che
    corrisponde a una data intensità dello stimolo fisico. È il processo
    attraverso cui la stimolazione dei recettori sensoriali produce impulsi
    neuronali collegati alle esperienze vissute dentro e fuori dal corpo.
\end{snippetdefinition}

\begin{snippetdefinition}{organizzazione-percettiva-definition}{Organizzazione percettiva}
    L'\emph{organizzazione percettiva} è la fase in cui il cervello integra i
    segnali raccolti dagli organi recettori con la conoscenza pregressa del
    mondo, formando una rappresentazione interna coerente dello stimolo esterno.
\end{snippetdefinition}

\begin{snippet}{identificazione-riconoscimento-oggetti}
    I processi di \emph{identificazione} e \emph{riconoscimento} attribuiscono
    significato ai percetti. Nel caso della visione, le prime fasi della
    percezione rispondono alla domanda su che aspetto abbia l'oggetto; le fasi
    successive rispondono invece a domande su che cosa sia l'oggetto e quale
    funzione abbia.
\end{snippet}

\begin{snippetdefinition}{stimolo-distale-definition}{Stimolo distale}
    Lo \emph{stimolo distale} è l'oggetto fisico reale collocato nell'ambiente,
    distante dall'osservatore.
\end{snippetdefinition}

\begin{snippetdefinition}{stimolo-prossimale-definition}{Stimolo prossimale}
    Lo \emph{stimolo prossimale} è la rappresentazione dello stimolo distale
    sugli organi sensoriali, per esempio l'immagine retinica nel caso della
    visione. È lo stimolo vicino all'osservatore da cui il sistema percettivo
    ricava direttamente le informazioni.
\end{snippetdefinition}

\begin{snippet}{stimolo-distale-prossimale}
    Il sistema percettivo mira a ricostruire lo stimolo distale, cioè l'oggetto
    reale presente nell'ambiente, a partire dallo stimolo prossimale. Nel caso
    della visione, questa ricostruzione è complessa perché l'immagine retinica è
    bidimensionale, mentre l'ambiente è tridimensionale.
\end{snippet}

\begin{snippetdefinition}{processi-bottom-up-definition}{Processi bottom-up}
    I \emph{processi bottom-up}, o processi dal basso verso l'alto, sono
    processi percettivi guidati dalle informazioni sensoriali provenienti dal
    mondo fisico. Sono detti anche processi basati sui dati, perché partono
    dall'evidenza sensoriale rilevata nel contesto.
\end{snippetdefinition}

\begin{snippetdefinition}{processi-top-down-definition}{Processi top-down}
    I \emph{processi top-down}, o processi dall'alto verso il basso, sono
    processi percettivi guidati da conoscenze, credenze, aspettative,
    motivazioni, obiettivi ed esperienze passate. Sono detti anche processi
    guidati dai concetti.
\end{snippetdefinition}

\begin{snippet}{bottom-up-top-down-percezione}
    La percezione deriva dall'interazione tra processi bottom-up e processi
    top-down. I primi portano al cervello le informazioni sensoriali raccolte
    dall'ambiente; i secondi usano conoscenze, aspettative e memoria per
    interpretare tali informazioni.
\end{snippet}

\includesnpt[width=58\%,src=/snippet/static-psychology/bottom-up-top-down-processing.jpg]{centered-img}

\end{document}
